% TOP_PROBLEM: {"id": 3286, "title": "Triangle-packing vs triangle edge-transversal.", "category_id": 3, "category": {"id": 3, "name": "graph_theory", "display_name": "Graph Theory", "description": "Problems involving graphs, networks, and their properties.", "slug": "graph-theory", "order_index": 3, "created_at": "2026-07-31T15:26:25.671Z"}, "status": "open", "problem_number": "OPG-46837", "set_id": 12}
% FIRST_POSTED: 2026-10-01
% ENTRY_KIND: statement_only
% UnsolvedMath ID 3286; source catalogue code OPG-46837
% !TeX program = lualatex
% Definition and sources only; this file is not an LLM solution attempt.
\documentclass[11pt,a4paper]{article}
\usepackage[margin=25mm]{geometry}
\usepackage{fontspec}
\setmainfont{lmroman10-regular.otf}[BoldFont=lmroman10-bold.otf,
 ItalicFont=lmroman10-italic.otf,BoldItalicFont=lmroman10-bolditalic.otf]
\usepackage{amsmath,amssymb,mathtools}
\usepackage{unicode-math}
\setmathfont{latinmodern-math.otf}
\AtBeginDocument{\let\setminus\smallsetminus}
\usepackage{xurl,hyperref}
\hypersetup{colorlinks=true,urlcolor=blue}
\setcounter{secnumdepth}{0}
\setlength{\parindent}{0pt}
\setlength{\parskip}{5pt}
\setlength{\emergencystretch}{3em}
\begin{document}
\section*{Triangle-packing vs triangle edge-transversal.}\label{3286}
{\small UnsolvedMath ID 3286 · OPG-46837 · Graph Theory · OpenGarden}
\subsection{Definitions and mathematical statement}
Let $G=(V,E)$ be a finite simple undirected graph. A triangle is the
three-edge subgraph on three pairwise adjacent vertices. A triangle
packing is a family of triangles with pairwise disjoint edge sets;
sharing vertices is permitted. Write $\nu_{\triangle}(G)$ for the largest
number of triangles in such a family.

A triangle edge-transversal is a set $F\subseteq E$ meeting the edge set
of every triangle. Equivalently, deleting $F$ leaves a triangle-free
graph. Write $\tau_{\triangle}(G)$ for the minimum size of a triangle
edge-transversal. Both parameters are zero when the graph has no
triangle.

Tuza's conjecture asks whether
\[
                    \tau_{\triangle}(G)\le2\nu_{\triangle}(G)
\]
for every finite simple graph $G$. Equivalently, if no packing has more
than $k$ triangles, some set of at most $2k$ edges destroys every
triangle. ``At most'' is the appropriate convention: graphs with fewer
than $2k$ edges need not admit a set of exactly that size.

The elementary inequalities are
$\nu_{\triangle}(G)\le\tau_{\triangle}(G)\le3\nu_{\triangle}(G)$.
For the upper bound, the union of the edges of a maximum packing
meets every triangle, since otherwise the packing could be enlarged.
The conjecture improves that factor three to two. For $K_4$, the
packing number is one and the transversal number is two, so the
proposed factor cannot be reduced. Both quantities are integral;
fractional assignments of weights to edges or triangles are different
relaxations of the problem.
\subsection{Short English statement}
Can all triangles be destroyed using at most twice as many edges as the maximum number of edge-disjoint triangles?
\subsection{Sources}
\begin{itemize}
\item[S1] ULAM.AI and the UnsolvedMath Contributors, supplied problems.json, record 3286 (OPG-46837), OpenGarden collection. Catalogue identity and source transcription; CC BY 4.0.
\url{https://huggingface.co/datasets/ulamai/UnsolvedMath}.
\item[S2] Open Problem Garden, Triangle-packing vs triangle edge-transversal.. Original problem and discussion; the mathematical formulation below is expanded from the supplied source record.
\url{http://www.openproblemgarden.org/op/triangle_packing_vs_triangle_edge_transversal}.
\end{itemize}
\end{document}
